\begin{proof}[Proof of \cref{obj:thm:adaptive-minimax}]
\emph{1. Uniform local approximation.}
We first establish the estimator bound uniformly over the prescribed adaptation interval. Set
\[
m:=\lceil\beta\rceil-1,\qquad
s:=\beta-m,\qquad
\mathcal A_{\mathrm{poly}}
:=\{\alpha\in\mathbb N_0^d:|\alpha|\le m\},
\qquad
p_{\mathrm{coef}}:=|\mathcal A_{\mathrm{poly}}|.
\]
Then \(0<s\le1\) and \(p_{\mathrm{coef}}>0\). By the identification and uniqueness conclusions of \cref{obj:lem:causal-identification}, every Hölder response witnessing membership of \(P\) in \(\mathcal P_\gamma\) agrees with \(\mu_{1,P}\) throughout \(\mathcal X\).

Here is the local approximation bound used below, including at cube boundaries. Define the affine map and transported response by
\[
T_{\mathrm{aff}}(x):=2x-\mathbf1
\quad(x\in\mathbb R^d),
\qquad
g_{\mathrm{aff}}(z)
:=\mu_{1,P}\!\left(\frac{z+\mathbf1}{2}\right)
\quad(z\in[-1,1]^d).
\]
Repeated coordinate differentiation of the affine composition gives, for every multiindex \(\alpha\) with \(|\alpha|\le m\),
\[
D^\alpha g_{\mathrm{aff}}(z)
=
2^{-|\alpha|}D^\alpha\mu_{1,P}\!\left(\frac{z+\mathbf1}{2}\right),
\qquad z\in[-1,1]^d.
\]
The chain rule proves this identity in the interior, and continuity of the derivative traces extends it to the boundary. Thus \cref{obj:ass:holder-derivative-bound} gives
\[
\sup_{z\in[-1,1]^d}|D^\alpha g_{\mathrm{aff}}(z)|
\le 2^{-|\alpha|}L\le L.
\]
For \(|\alpha|=m\) and \(z,z'\in[-1,1]^d\), the identity and \cref{obj:ass:holder-modulus} give
\[
\begin{aligned}
|D^\alpha g_{\mathrm{aff}}(z)-D^\alpha g_{\mathrm{aff}}(z')|
&=2^{-m}\left|
D^\alpha\mu_{1,P}\!\left(\frac{z+\mathbf1}{2}\right)
-D^\alpha\mu_{1,P}\!\left(\frac{z'+\mathbf1}{2}\right)
\right|
\\
&\le 2^{-m}L
\left\|\frac{z-z'}2\right\|_\infty^s
\\
&\le L\|z-z'\|_\infty^s.
\end{aligned}
\]
Affine composition also preserves the continuous derivative traces through order \(m\). Consequently, defining
\[
K_{\mathrm{aff}}:=1,
\]
we obtain intrinsic Hölder radius \(K_{\mathrm{aff}}L\) for \(g_{\mathrm{aff}}\) on \([-1,1]^d\). \Cref{obj:lem:global-holder-extension}, applied with derivative order \(m\), exponent \(s\), and radius \(K_{\mathrm{aff}}L\), supplies a constant \(K_{\mathrm{ext}}>0\) and a function \(g_{\mathrm{ext}}\in C^m(\mathbb R^d)\) satisfying
\[
g_{\mathrm{ext}}(z)=g_{\mathrm{aff}}(z)
\quad(z\in[-1,1]^d),
\]
\[
\max_{0\le k_{\mathrm{der}}\le m}
\sup_{z\in\mathbb R^d}
\|\mathrm D^{k_{\mathrm{der}}}g_{\mathrm{ext}}(z)\|_{\mathrm{op}}
\le K_{\mathrm{ext}}K_{\mathrm{aff}}L,
\]
and
\[
\|\mathrm D^m g_{\mathrm{ext}}(z)
-\mathrm D^m g_{\mathrm{ext}}(z')\|_{\mathrm{op}}
\le
K_{\mathrm{ext}}K_{\mathrm{aff}}L
\|z-z'\|_\infty^s
\quad(z,z'\in\mathbb R^d).
\]
Here \(\mathrm D^{k_{\mathrm{der}}}\) denotes the Fréchet derivative of the indicated order, with order zero interpreted as the function itself.

For a partition cube \(Q\) of mesh \(0<h\le1\), define
\[
\vartheta_{Q,\alpha}
:=
\frac{(2h)^{|\alpha|}}{\alpha!}
D^\alpha g_{\mathrm{ext}}(T_{\mathrm{aff}}(a_Q)),
\qquad \alpha\in\mathcal A_{\mathrm{poly}},
\]
and
\[
\Pi_Q(x)
:=
U\!\left(\frac{x-a_Q}{h}\right)^\top\vartheta_Q,
\qquad x\in Q.
\]
We derive the remainder bound along the expansion segment. Since \(\beta>1\), we have \(m\ge1\). For \(x\in Q\), define
\[
z_{Q,0}:=T_{\mathrm{aff}}(a_Q),
\qquad
v_{Q,x}:=2(x-a_Q),
\qquad
\varphi_{Q,x}(t_{\mathrm{seg}})
:=g_{\mathrm{ext}}(z_{Q,0}+t_{\mathrm{seg}}v_{Q,x})
\quad(t_{\mathrm{seg}}\in\mathbb R).
\]
The chain rule gives
\[
\varphi_{Q,x}^{(k_{\mathrm{der}})}(t_{\mathrm{seg}})
=
\mathrm D^{k_{\mathrm{der}}}g_{\mathrm{ext}}
(z_{Q,0}+t_{\mathrm{seg}}v_{Q,x})
[\underbrace{v_{Q,x},\ldots,v_{Q,x}}_{k_{\mathrm{der}}\text{ arguments}}],
\qquad
0\le k_{\mathrm{der}}\le m,
\quad t_{\mathrm{seg}}\in\mathbb R,
\]
with order zero interpreted as the function value. Expanding the symmetric derivatives in coordinate directions and using the definition of \(\vartheta_Q\) yields
\[
\sum_{k_{\mathrm{der}}=0}^{m}
\frac{\varphi_{Q,x}^{(k_{\mathrm{der}})}(0)}{k_{\mathrm{der}}!}
=
\sum_{\alpha\in\mathcal A_{\mathrm{poly}}}
\frac{D^\alpha g_{\mathrm{ext}}(z_{Q,0})}{\alpha!}
\prod_{i=1}^{d}(v_{Q,x,i})^{\alpha_i}
=\Pi_Q(x),
\qquad
\varphi_{Q,x}(1)=\mu_{1,P}(x).
\]
Taylor's theorem with the order-\(m\) Lagrange remainder gives a point satisfying
\[
\exists\,t_{Q,x}\in(0,1):\qquad
\varphi_{Q,x}(1)
-
\sum_{k_{\mathrm{der}}=0}^{m-1}
\frac{\varphi_{Q,x}^{(k_{\mathrm{der}})}(0)}{k_{\mathrm{der}}!}
=
\frac{\varphi_{Q,x}^{(m)}(t_{Q,x})}{m!}.
\]
Choose such a point and define the degree-\(m\) remainder by
\[
\begin{aligned}
R_{\mathrm{Taylor},Q,x}
&:=\mu_{1,P}(x)-\Pi_Q(x)\\
&=\frac{\varphi_{Q,x}^{(m)}(t_{Q,x})
-\varphi_{Q,x}^{(m)}(0)}{m!}\\
&=\frac1{m!}
\left(
\mathrm D^m g_{\mathrm{ext}}(z_{Q,0}+t_{Q,x}v_{Q,x})
-\mathrm D^m g_{\mathrm{ext}}(z_{Q,0})
\right)
[\underbrace{v_{Q,x},\ldots,v_{Q,x}}_{m\text{ arguments}}].
\end{aligned}
\]
The displayed Hölder modulus for the extended function therefore implies
\[
\begin{aligned}
|R_{\mathrm{Taylor},Q,x}|
&\le
\frac{K_{\mathrm{ext}}K_{\mathrm{aff}}L}{m!}
\|t_{Q,x}v_{Q,x}\|_\infty^s
\|v_{Q,x}\|_\infty^m\\
&=
\frac{K_{\mathrm{ext}}K_{\mathrm{aff}}L}{m!}
 t_{Q,x}^{s}\|v_{Q,x}\|_\infty^{m+s}\\
&\le
\frac{K_{\mathrm{ext}}K_{\mathrm{aff}}L}{m!}
\|v_{Q,x}\|_\infty^\beta,
\end{aligned}
\]
since \(0<t_{Q,x}<1\), \(s>0\), and \(m+s=\beta\). This bound also applies when \(v_{Q,x}=0\). Define
\[
C_{\mathrm{Taylor}}:=\frac1{m!}\ge0.
\]
Because \(\|v_{Q,x}\|_\infty\le2h\), we obtain, uniformly in the expansion centre,
\[
\begin{aligned}
|\mu_{1,P}(x)-\Pi_Q(x)|
&=
\left|
g_{\mathrm{ext}}\!\left(
T_{\mathrm{aff}}(a_Q)+2h\frac{x-a_Q}{h}
\right)
-
U\!\left(\frac{x-a_Q}{h}\right)^\top\vartheta_Q
\right|
\\
&\le
C_{\mathrm{Taylor}}K_{\mathrm{ext}}K_{\mathrm{aff}}L(2h)^\beta
=
C_{\mathrm{bias}}Lh^\beta,
\end{aligned}
\]
where
\[
C_{\mathrm{bias}}
:=
C_{\mathrm{Taylor}}K_{\mathrm{ext}}K_{\mathrm{aff}}2^\beta.
\]
The normalized coordinates belong to \([0,1]^d\), and the affine image of \(x\) belongs to \([-1,1]^d\), which justify the equality and the approximation. This construction also covers integer \(\beta\), for which \(s=1\).
% lean: CausalSmith.Stat.WeakOverlap.exists_uniform_holderBall_local_monoVec_approx

\emph{2. Conditional bias and supremum loss.}
Apply \cref{obj:lem:selected-mesh}. Denote its uniform mesh and coefficient constants by \(K_{\mathrm{mesh}},K_{\mathrm{coef}}>0\), and its sample-size threshold by \(n_{\mathrm{env}}\). Enlarge these constants and the threshold as required by the constructions below, taking \(n_{\mathrm{env}}\ge2\). Define
\[
n_0:=\max\{n_{\mathrm{env}},1\},
\qquad
\mathcal G_n:=\sigma((X_i,A_i)_{i=1}^n).
\]
For \(n\ge n_0\), \(\gamma\in\Gamma\), and \(P\in\mathcal P_\gamma\), choose the simultaneous treated-count event underlying the mesh and coefficient bounds as follows. For all candidate microcells, define
\[
q_{Q\ell}:=P\{A=1,\ X\in E_{Q\ell}\},
\qquad h\in\mathcal H_n,\quad Q\in\mathcal Q_h,\quad 1\le\ell\le J,
\]
\[
M_{\mathrm{count},n}
:=\sum_{h\in\mathcal H_n}J|\mathcal Q_h|,
\qquad
s_{\mathrm{count},n}
:=\log\!\left(2(M_{\mathrm{count},n}+1)n^2\right)>0,
\]
and take
\[
\mathcal E_n
:=
\bigcap_{h\in\mathcal H_n}
\bigcap_{Q\in\mathcal Q_h}
\bigcap_{\ell=1}^J
\left\{
\left|N_{Q\ell}-nq_{Q\ell}\right|
<\frac{nq_{Q\ell}}2+4s_{\mathrm{count},n}
\right\}.
\]
The Bernstein bound for each treated count and the finite union bound give
\[
\begin{aligned}
P^{\otimes n}(\mathcal E_n^c)
&\le 2M_{\mathrm{count},n}e^{-s_{\mathrm{count},n}}
\\
&=\frac{M_{\mathrm{count},n}}{M_{\mathrm{count},n}+1}n^{-2}
\le n^{-2}.
\end{aligned}
\]
We establish the selected-mesh and treated-count bounds on this particular intersection. Let \(\eta\) be the reference microcube width in \cref{obj:def:template}. Each microcell at mesh \(h\) has volume \((\eta h)^d\), so \cref{obj:ass:density} and the measurable-set bound in \cref{obj:lem:ordered-mass} give
\[
\begin{aligned}
P(X\in E_{Q\ell})&\ge c_f(\eta h)^d,\\
q_{Q\ell}
&\ge \frac{\gamma-1}{\gamma}C^{-1/(\gamma-1)}
\bigl(c_f(\eta h)^d\bigr)^{\gamma/(\gamma-1)}
=\kappa_{\mathrm{cell}}(\gamma)h^{D_\gamma},
\end{aligned}
\]
where
\[
\kappa_{\mathrm{cell}}(\gamma)
:=\frac{\gamma-1}{\gamma}C^{-1/(\gamma-1)}
(c_f\eta^d)^{\gamma/(\gamma-1)},
\qquad
\kappa_{\mathrm{cell},\Gamma}
:=\min_{\gamma'\in\Gamma}\kappa_{\mathrm{cell}}(\gamma')>0.
\]
The positivity of the minimum follows from continuity and positivity of the displayed coefficient on the compact interval \(\Gamma\). Choose a fixed enlargement constant satisfying
\[
A_{\mathrm{or}}\ge1,
\qquad
4\le\kappa_{\mathrm{cell},\Gamma}
A_{\mathrm{or}}^{2\beta+D_{\gamma_{\max}}}.
\]
Such a choice exists because \(2\beta+D_{\gamma_{\max}}>0\). Since
\[
D_\gamma=d+\frac d{\gamma-1},
\qquad
D_{\gamma_{\max}}\le D_\gamma\le D_{\gamma_{\min}},
\qquad
h_{\gamma_{\max},n}\le h_{\gamma,n}\le h_{\gamma_{\min},n},
\]
we have
\[
4\le\kappa_{\mathrm{cell},\Gamma}
A_{\mathrm{or}}^{2\beta+D_\gamma}
\quad(\gamma\in\Gamma).
\]

The candidate collection in \cref{obj:synth_3} has largest index
\[
j_{\mathrm{cand},n}
:=\left\lfloor\frac{\log_2 n}{2\beta+d}\right\rfloor,
\qquad
\mathcal H_n=\{2^{-j}:0\le j\le j_{\mathrm{cand},n}\}.
\]
The floor inequality gives
\[
2^{-j_{\mathrm{cand},n}}
\le2n^{-1/(2\beta+d)}.
\]
Because \(d>0\) and \(\gamma_{\max}>1\),
\[
D_{\gamma_{\max}}>d,
\qquad
\frac{2n^{-1/(2\beta+d)}}{A_{\mathrm{or}}h_{\gamma_{\max},n}}
=\frac2{A_{\mathrm{or}}}
 n^{-1/(2\beta+d)+1/(2\beta+D_{\gamma_{\max}})}
\longrightarrow0,
\qquad
A_{\mathrm{or}}h_{\gamma_{\min},n}\longrightarrow0.
\]
After enlarging \(n_{\mathrm{env}}\), these relations ensure, uniformly over \(\gamma\in\Gamma\), that
\[
2^{-j_{\mathrm{cand},n}}
\le A_{\mathrm{or}}h_{\gamma,n}\le1.
\]
Define an oracle candidate by
\[
j_{\mathrm{or}}(n,\gamma)
:=\max\{j\in\{0,\ldots,j_{\mathrm{cand},n}\}:
2^{-j}\ge A_{\mathrm{or}}h_{\gamma,n}\},
\qquad
h_{\mathrm{or},n,\gamma}:=2^{-j_{\mathrm{or}}(n,\gamma)}.
\]
The set defining the maximum contains zero. If its maximum is smaller than \(j_{\mathrm{cand},n}\), the next dyadic width is smaller than \(A_{\mathrm{or}}h_{\gamma,n}\); if the maximum equals \(j_{\mathrm{cand},n}\), the preceding lower bound on the finest width forces equality with the target width. Thus
\[
h_{\mathrm{or},n,\gamma}\in\mathcal H_n,
\qquad
A_{\mathrm{or}}h_{\gamma,n}
\le h_{\mathrm{or},n,\gamma}
\le2A_{\mathrm{or}}h_{\gamma,n}.
\]

We next verify the logarithmic budget. The template and dyadic partition have
\[
J=(m+1)^d,
\qquad
|\mathcal Q_{2^{-j}}|=2^{jd},
\qquad
M_{\mathrm{count},n}
=J\sum_{j=0}^{j_{\mathrm{cand},n}}2^{jd}.
\]
Since \(2\beta+d\ge1\), the definition of \(j_{\mathrm{cand},n}\) gives
\[
j_{\mathrm{cand},n}\le2^{j_{\mathrm{cand},n}}\le n,
\qquad
M_{\mathrm{count},n}\le(n+1)n^dJ.
\]
For \(n\ge2\), using \(n+1\le2n\) and \(n^{d+1}\ge1\), we obtain
\[
2(M_{\mathrm{count},n}+1)n^2
\le8(J+1)n^{d+3},
\qquad
s_{\mathrm{count},n}
\le(d+3)\log n+\log(8(J+1)).
\]
Define
\[
U_{\mathrm{count}}:=24(d+3),
\qquad
V_{\mathrm{count}}:=24\log(8(J+1)),
\qquad
\rho_{\mathrm{budget}}
:=\frac{2\beta}{2\beta+D_{\gamma_{\min}}}>0.
\]
The elementary bound
\[
\log n\le\frac2{\rho_{\mathrm{budget}}}
 n^{\rho_{\mathrm{budget}}/2}
\]
and divergence of positive powers allow a further uniform enlargement of \(n_{\mathrm{env}}\) such that
\[
2U_{\mathrm{count}}\frac2{\rho_{\mathrm{budget}}}
\le(2A_{\mathrm{or}})^{-2\beta}
 n^{\rho_{\mathrm{budget}}/2},
\qquad
2V_{\mathrm{count}}
\le(2A_{\mathrm{or}})^{-2\beta}n^{\rho_{\mathrm{budget}}}.
\]
Consequently,
\[
\begin{aligned}
24s_{\mathrm{count},n}
&\le U_{\mathrm{count}}\log n+V_{\mathrm{count}}\\
&\le(2A_{\mathrm{or}})^{-2\beta}n^{\rho_{\mathrm{budget}}}
=(2A_{\mathrm{or}}h_{\gamma_{\min},n})^{-2\beta}\\
&\le(2A_{\mathrm{or}}h_{\gamma,n})^{-2\beta}
\le h_{\mathrm{or},n,\gamma}^{-2\beta}.
\end{aligned}
\]
The population-mass bound also gives, for every microcell of this oracle candidate,
\[
\begin{aligned}
nq_{Q\ell}
&\ge n\kappa_{\mathrm{cell},\Gamma}
 h_{\mathrm{or},n,\gamma}^{D_\gamma}\\
&=\kappa_{\mathrm{cell},\Gamma}
 n h_{\mathrm{or},n,\gamma}^{2\beta+D_\gamma}
 h_{\mathrm{or},n,\gamma}^{-2\beta}\\
&\ge\kappa_{\mathrm{cell},\Gamma}
 A_{\mathrm{or}}^{2\beta+D_\gamma}
 h_{\mathrm{or},n,\gamma}^{-2\beta}
\ge4h_{\mathrm{or},n,\gamma}^{-2\beta},
\end{aligned}
\]
where \(nh_{\gamma,n}^{2\beta+D_\gamma}=1\).

On \(\mathcal E_n\), the defining deviation inequality therefore yields, at the oracle candidate,
\[
N_{Q\ell}
>\frac{nq_{Q\ell}}2-4s_{\mathrm{count},n}
\ge2h_{\mathrm{or},n,\gamma}^{-2\beta}
 -4s_{\mathrm{count},n}
\ge h_{\mathrm{or},n,\gamma}^{-2\beta}.
\]
Thus this candidate is feasible. The smallest-feasible-mesh rule of \cref{obj:def:mesh-selector} gives
\[
0<\widehat h\le h_{\mathrm{or},n,\gamma}
\le2A_{\mathrm{or}}h_{\gamma,n},
\qquad
24s_{\mathrm{count},n}
\le h_{\mathrm{or},n,\gamma}^{-2\beta}
\le\widehat h^{-2\beta}.
\]
For a microcell of the selected mesh, feasibility and the same deviation inequality imply
\[
N_{Q\ell}\ge\widehat h^{-2\beta}
\ge24s_{\mathrm{count},n},
\qquad
\frac{nq_{Q\ell}}2
<N_{Q\ell}+4s_{\mathrm{count},n}
\le\frac76N_{Q\ell}.
\]
In particular,
\[
N_{Q\ell}\ge\frac{nq_{Q\ell}}5
\quad(Q\in\mathcal Q_{\widehat h},\ 1\le\ell\le J).
\]
Taking \(K_{\mathrm{mesh}}\ge2A_{\mathrm{or}}\), we have established on the displayed count intersection that
\[
0<\widehat h\le K_{\mathrm{mesh}}h_{\gamma,n},
\qquad
N_{Q\ell}\ge\frac{nq_{Q\ell}}5.
\]
% lean: CausalSmith.Stat.WeakOverlap.selectedCountEvent_uniform_design_saturated
% lean: CausalSmith.Stat.WeakOverlap.selectedCount_deterministic_bridge
Every count in the defining finite intersection is measurable with respect to \(\mathcal G_n\), and every population mass is fixed under \(P\). Hence \(\mathcal E_n\in\mathcal G_n\). Conditioning therefore fixes the sampled design, the selected mesh, its treated counts, and membership in \(\mathcal E_n\).

First, the response bound holds everywhere on the cube. Define
\[
g_{\mathrm{clip}}(x)
:=
\operatorname{clip}_{[-B,B]}(\mu_{1,P}(x)),
\qquad x\in\mathcal X.
\]
The conditional-mean bound in \cref{obj:ass:bounded-potential-outcomes} and identification give
\[
g_{\mathrm{clip}}(x)
=
\mu_{1,P}(x)
=
\mathbb E_P[Y\mid A=1,X=x]
\quad\text{for }P_X\text{-almost every }x.
\]
Both functions are continuous on the cube. The uniqueness conclusion of \cref{obj:lem:causal-identification} therefore gives
\[
g_{\mathrm{clip}}(x)=\mu_{1,P}(x),
\qquad
|\mu_{1,P}(x)|\le B
\quad(x\in\mathcal X).
\]

On \(\mathcal E_n\), set
\[
h:=\widehat h,
\qquad
\bar c_Q:=\mathbb E[\widehat c_Q\mid\mathcal G_n],
\qquad Q\in\mathcal Q_h.
\]
At each fixed design, this conditional expectation evaluates the coefficients at that design's selected mesh. Feasibility gives positive treated counts. For each treated observation used in the fit on \(Q\), define
\[
U_{i,Q}:=U\!\left(\frac{X_i-a_Q}{h}\right),
\qquad
R_{\mathrm{res},i,Q}
:=\mu_{1,P}(X_i)-U_{i,Q}^{\top}\vartheta_Q.
\]
The local polynomial approximation gives
\[
|R_{\mathrm{res},i,Q}|\le C_{\mathrm{bias}}Lh^\beta.
\]
Identification and the equal-cell normal equations yield, almost surely,
\[
\bar c_Q-\vartheta_Q
=
\widehat G_Q^{-1}
\frac1J\sum_{\ell=1}^J\frac1{N_{Q\ell}}
\sum_{\substack{i:A_i=1\\X_i\in E_{Q\ell}}}
U_{i,Q}R_{\mathrm{res},i,Q}.
\]
Every monomial coordinate has absolute value at most one on the normalized cube, so
\[
\|U_{i,Q}\|_2\le\sqrt{p_{\mathrm{coef}}}.
\]
Consequently, \cref{obj:lem:template-conditioning} implies
\[
\|\bar c_Q-\vartheta_Q\|_\infty
\le
\|\bar c_Q-\vartheta_Q\|_2
\le
\frac{2\sqrt{p_{\mathrm{coef}}}}{\lambda_0}
C_{\mathrm{bias}}Lh^\beta.
\]
Define the fixed constants
\[
K_{\mathrm{inv}}
:=\frac{2\sqrt{p_{\mathrm{coef}}}}{\lambda_0},
\qquad
K_{\mathrm{bias}}
:=p_{\mathrm{coef}}K_{\mathrm{inv}}C_{\mathrm{bias}}L
+C_{\mathrm{bias}}L.
\]

For the selected design, define the maximum coefficient fluctuation by
\[
S_{\mathrm{coef}}
:=
\max_{Q\in\mathcal Q_h}
\|\widehat c_Q-\bar c_Q\|_\infty.
\]
This maximum uses the selected mesh fixed by the conditioning design. Since
\[
|U(z)^\top(\widehat c_Q-\vartheta_Q)|
\le
p_{\mathrm{coef}}
\bigl(S_{\mathrm{coef}}+K_{\mathrm{inv}}C_{\mathrm{bias}}Lh^\beta\bigr),
\qquad z\in[0,1]^d,
\]
and clipping contracts distance to every value in \([-B,B]\), we obtain
\[
\sup_{x\in\mathcal X}
|\widehat\mu_n(x)-\mu_{1,P}(x)|
\le
p_{\mathrm{coef}}S_{\mathrm{coef}}+K_{\mathrm{bias}}h^\beta.
\]
We derive the conditional coefficient envelope on this same count event. For the selected design, define
\[
q_Q:=\min_{1\le\ell\le J}q_{Q\ell},
\qquad
\rho_\gamma:=\frac1{\gamma-1},
\]
and order the selected cubes so that
\[
\{Q_{(k)}:1\le k\le h^{-d}\}=\mathcal Q_h,
\qquad
q_{(k)}:=q_{Q_{(k)}},
\qquad
q_{(1)}\le\cdots\le q_{(h^{-d})}.
\]
Here \(h^{-d}=|\mathcal Q_h|\ge1\). The ordered-mass conclusion of \cref{obj:lem:ordered-mass} supplies a fixed constant satisfying
\[
\kappa_\Gamma>0,
\qquad
q_{(k)}\ge\kappa_\Gamma h^{D_\gamma}k^{\rho_\gamma}
\quad(1\le k\le h^{-d}).
\]

Define the treated residuals and coefficient weights by
\[
\epsilon_{\mathrm{tr},i}
:=\mathbf1_{\{A_i=1\}}\bigl(Y_i-\mu_{1,P}(X_i)\bigr),
\]
\[
w_{\mathrm{fit},i,Q,\alpha}
:=\frac1J\sum_{\ell=1}^J
\frac{\mathbf1_{\{A_i=1,\ X_i\in E_{Q\ell}\}}}{N_{Q\ell}}
\bigl(\widehat G_Q^{-1}U_{i,Q}\bigr)_\alpha,
\qquad \alpha\in\mathcal A_{\mathrm{poly}}.
\]
These weights are evaluated only at the positive feasible mesh on \(\mathcal E_n\), where every denominator is positive. Conditional on \(\mathcal G_n\), the residuals are independent by iid sampling, and identification together with \cref{obj:ass:bounded-potential-outcomes} gives, almost surely,
\[
\mathbb E[\epsilon_{\mathrm{tr},i}\mid\mathcal G_n]=0,
\qquad
\mathbb E[e^{t\epsilon_{\mathrm{tr},i}}\mid\mathcal G_n]
\le e^{B^2t^2/2}
\quad(t\in\mathbb R).
\]
For an untreated observation the residual is zero. The equal-cell coefficient formula consequently gives
\[
\widehat c_{Q,\alpha}-\bar c_{Q,\alpha}
=\sum_{i=1}^n w_{\mathrm{fit},i,Q,\alpha}
\epsilon_{\mathrm{tr},i}.
\]
The inverse bound in \cref{obj:lem:template-conditioning} and \(\|U_{i,Q}\|_2\le\sqrt{p_{\mathrm{coef}}}\) imply
\[
\left|\frac1J
(\widehat G_Q^{-1}U_{i,Q})_\alpha\right|
\le\frac{K_{\mathrm{inv}}}{J}.
\]
Cauchy--Schwarz over the \(J\) microcells, followed by summation over their treated observations, gives
\[
\begin{aligned}
\sum_{i=1}^n w_{\mathrm{fit},i,Q,\alpha}^2
&\le J\left(\frac{K_{\mathrm{inv}}}{J}\right)^2
\sum_{\ell=1}^J\frac1{N_{Q\ell}}\\
&\le K_{\mathrm{inv}}^2
\min\left(h^{2\beta},\frac5{nq_Q}\right).
\end{aligned}
\]
The last line uses both \(N_{Q\ell}\ge h^{-2\beta}\) and the established bound \(N_{Q\ell}\ge nq_{Q\ell}/5\). Multiplying the conditional exponential moments of the independent weighted residuals yields
\[
\mathbb E\!\left[
 e^{t(\widehat c_{Q,\alpha}-\bar c_{Q,\alpha})}
 \,\middle|\,\mathcal G_n\right]
\le\exp\!\left\{
\frac{B^2t^2}{2}\sum_{i=1}^n
w_{\mathrm{fit},i,Q,\alpha}^2\right\}.
\]
Define
\[
K_{\mathrm{mgf}}
:=B^2K_{\mathrm{inv}}^2\max\{1,5/\kappa_\Gamma\}>0.
\]
The ordered-mass bound and
\(\min\{u,cv\}\le\max\{1,c\}\min\{u,v\}\) for nonnegative \(u,v,c\) now give
\[
\mathbb E\!\left[
 e^{t(\widehat c_{Q_{(k)},\alpha}-\bar c_{Q_{(k)},\alpha})}
 \,\middle|\,\mathcal G_n\right]
\le\exp\!\left\{
\frac{K_{\mathrm{mgf}}t^2}{2}
\min\left(h^{2\beta},n^{-1}h^{-D_\gamma}k^{-\rho_\gamma}\right)
\right\}.
\]
All these exponential moments are finite; their two-sided finiteness also gives conditional integrability of the finite coefficient maxima.
% lean: CausalSmith.Stat.WeakOverlap.selectedMesh_saturated_rank_mgf

To apply the maximal inequality in \cref{obj:lem:selected-mesh} with a constant uniform in \(\gamma\), define
\[
\rho_{\min}:=\frac1{\gamma_{\max}-1},
\qquad
\rho_{\max}:=\frac1{\gamma_{\min}-1},
\qquad
R_{\mathrm{max}}:=\frac{\rho_{\max}}{\rho_{\min}},
\]
\[
a_{\mathrm{max}}:=\sqrt{K_{\mathrm{mgf}}}\,h^\beta,
\qquad
b_{\mathrm{max}}:=\sqrt{K_{\mathrm{mgf}}}\,n^{-1/2}h^{-D_\gamma/2}.
\]
Then
\[
0<\rho_{\min}\le\rho_\gamma\le\rho_{\max},
\qquad R_{\mathrm{max}}\ge1,
\qquad
 a_{\mathrm{max}}^2=K_{\mathrm{mgf}}h^{2\beta},
\qquad
 b_{\mathrm{max}}^2=K_{\mathrm{mgf}}n^{-1}h^{-D_\gamma}.
\]
Since \(k^{-\rho_\gamma}\le k^{-\rho_{\min}}\) for \(k\ge1\), the same moment bound satisfies the maximal inequality at the fixed exponent \(\rho_{\min}\). Let \(K_{\min,\mathrm{max}}>0\) be its constant, and define
\[
M_{\mathrm{max}}:=K_{\min,\mathrm{max}}\sqrt{R_{\mathrm{max}}}.
\]
The logarithmic factors satisfy
\[
\begin{aligned}
1+\max\left\{0,\log\left(
 (b_{\mathrm{max}}/a_{\mathrm{max}})^{2/\rho_{\min}}
\right)\right\}
\le R_{\mathrm{max}}\left[
1+\max\left\{0,\log\left(
 (b_{\mathrm{max}}/a_{\mathrm{max}})^{2/\rho_\gamma}
\right)\right\}\right].
\end{aligned}
\]
Indeed, when \(b_{\mathrm{max}}/a_{\mathrm{max}}\ge1\), expand the logarithms and use \(\rho_\gamma\le\rho_{\max}\); when this ratio is smaller than one, both positive-part logarithms vanish. Applying the scalar maximal inequality to each coefficient and using
\[
S_{\mathrm{coef}}
\le\sum_{\alpha\in\mathcal A_{\mathrm{poly}}}
\max_{Q\in\mathcal Q_h}
|\widehat c_{Q,\alpha}-\bar c_{Q,\alpha}|,
\]
we obtain
\[
\mathbb E[S_{\mathrm{coef}}\mid\mathcal G_n]
\le p_{\mathrm{coef}}M_{\mathrm{max}}a_{\mathrm{max}}
\sqrt{1+\max\left\{0,\log\left(
(b_{\mathrm{max}}/a_{\mathrm{max}})^{2/\rho_\gamma}
\right)\right\}}.
\]

The oracle-mesh identity gives
\[
n^{-1/2}=h_{\gamma,n}^{\beta+D_\gamma/2},
\qquad
\frac{b_{\mathrm{max}}}{a_{\mathrm{max}}}
=\left(\frac{h_{\gamma,n}}h\right)^{\beta+D_\gamma/2}.
\]
Define
\[
\zeta_{\mathrm{log},\gamma}:=2\beta(\gamma-1)+d\gamma,
\qquad
\zeta_{\mathrm{log},\Gamma}
:=\max\{1,2\beta(\gamma_{\max}-1)+d\gamma_{\max}\}.
\]
Then
\[
\log\left((b_{\mathrm{max}}/a_{\mathrm{max}})^{2/\rho_\gamma}\right)
=\zeta_{\mathrm{log},\gamma}\log(h_{\gamma,n}/h),
\qquad
0<\zeta_{\mathrm{log},\gamma}\le\zeta_{\mathrm{log},\Gamma},
\]
and hence
\[
\mathbb E[S_{\mathrm{coef}}\mid\mathcal G_n]
\le p_{\mathrm{coef}}M_{\mathrm{max}}
\sqrt{K_{\mathrm{mgf}}\zeta_{\mathrm{log},\Gamma}}\,
 h^\beta\sqrt{1+\max\{0,\log(h_{\gamma,n}/h)\}}.
\]
To absorb the remaining logarithm, define
\[
z_{\mathrm{mesh}}:=h/h_{\gamma,n},
\qquad
0<z_{\mathrm{mesh}}\le K_{\mathrm{mesh}}.
\]
If \(z_{\mathrm{mesh}}\le1\), then
\[
\log(1/z_{\mathrm{mesh}})\le1/z_{\mathrm{mesh}}-1,
\qquad
\sqrt{1+\log(1/z_{\mathrm{mesh}})}\le1/z_{\mathrm{mesh}},
\qquad
z_{\mathrm{mesh}}^\beta\le z_{\mathrm{mesh}},
\]
so their product is at most one. If \(z_{\mathrm{mesh}}>1\), the positive-part logarithm vanishes and \(z_{\mathrm{mesh}}^\beta\le K_{\mathrm{mesh}}^\beta\). Thus
\[
h^\beta\sqrt{1+\max\{0,\log(h_{\gamma,n}/h)\}}
\le\max\{1,K_{\mathrm{mesh}}^\beta\}r_{\gamma,n}.
\]
Choose the initially enlarged coefficient constant to satisfy
\[
K_{\mathrm{coef}}\ge
p_{\mathrm{coef}}M_{\mathrm{max}}
\sqrt{K_{\mathrm{mgf}}\zeta_{\mathrm{log},\Gamma}}
\max\{1,K_{\mathrm{mesh}}^\beta\}.
\]
We have therefore established, almost surely on the explicitly chosen count event,
\[
\mathbb E[S_{\mathrm{coef}}\mid\mathcal G_n]
\le K_{\mathrm{coef}}r_{\gamma,n}
\quad\text{on }\mathcal E_n.
\]
% lean: CausalSmith.Stat.WeakOverlap.selectedMesh_saturated_conditional_envelope
Also,
\[
h^\beta
\le K_{\mathrm{mesh}}^\beta r_{\gamma,n}
\le \max\{1,K_{\mathrm{mesh}}^\beta\}r_{\gamma,n}.
\]
Thus, defining
\[
K_{\mathrm{cond}}
:=
p_{\mathrm{coef}}K_{\mathrm{coef}}
+
K_{\mathrm{bias}}\max\{1,K_{\mathrm{mesh}}^\beta\}
+1,
\]
we have \(K_{\mathrm{cond}}>0\) and
\[
\mathbb E\!\left[
\sup_{x\in\mathcal X}
|\widehat\mu_n(x)-\mu_{1,P}(x)|
\,\middle|\,\mathcal G_n
\right]
\le K_{\mathrm{cond}}r_{\gamma,n}
\quad\text{almost surely on }\mathcal E_n.
\]
The finite mesh selection and polynomial fits are measurable, and the conditional integrability needed here is supplied by \cref{obj:lem:selected-mesh}.
% lean: CausalSmith.Stat.WeakOverlap.selectedMesh_saturated_conditional_cubeLoss

\emph{3. Uniform expected risk.}
For every sample, including samples with \(\widehat h=0\), the estimator in \cref{obj:def:estimator} takes values in \([-B,B]\). Define, on the entire sample space,
\[
Z_{\mathrm{loss}}
:=
\sup_{x\in\mathcal X}
|\widehat\mu_n(x)-\mu_{1,P}(x)|,
\qquad
F_{\mathrm{loss}}
:=
\mathbb E[Z_{\mathrm{loss}}\mid\mathcal G_n].
\]
The preceding response bound gives
\[
0\le Z_{\mathrm{loss}}\le2B,
\qquad
0\le F_{\mathrm{loss}}\le2B.
\]
The supremum loss is measurable and bounded. The tower property and the bounds on the event and its complement therefore give
\[
\begin{aligned}
\mathbb E_{P^{\otimes n}}Z_{\mathrm{loss}}
&=\mathbb E_{P^{\otimes n}}F_{\mathrm{loss}}
\\
&\le
K_{\mathrm{cond}}r_{\gamma,n}
P^{\otimes n}(\mathcal E_n)
+
2B P^{\otimes n}(\mathcal E_n^c)
\\
&\le K_{\mathrm{cond}}r_{\gamma,n}+2Bn^{-2}.
\end{aligned}
\]
For every \(n\ge1\), \(\beta>0\), and \(\gamma>1\),
\[
D_\gamma\ge0,
\qquad
2\beta+D_\gamma>0,
\qquad
\frac{\beta}{2\beta+D_\gamma}\le1,
\]
so monotonicity of powers with base \(n\ge1\) gives
\[
n^{-2}
\le n^{-\beta/(2\beta+D_\gamma)}
=r_{\gamma,n}.
\]
Set
\[
K:=K_{\mathrm{cond}}+2B.
\]
Then \(K>0\), and
\[
\mathbb E_{P^{\otimes n}}
\left[
\sup_{x\in\mathcal X}
|\widehat\mu_n(x)-\mu_{1,P}(x)|
\right]
\le Kr_{\gamma,n}
\]
uniformly for \(n\ge n_0\), \(\gamma\in\Gamma\), and \(P\in\mathcal P_\gamma\). The count rule and clipped fit in \cref{obj:def:mesh-selector,obj:def:estimator} use the stated inputs \(n,d,\beta,B,\mathcal D_n\).
% lean: CausalSmith.Stat.WeakOverlap.estimatorSupRisk_uniform_upper

\emph{4. The bounded-subclass lower bound.}
Fix \(\gamma>1\) and \(x_0\in\mathcal X\). Define the subclass and its pointwise minimax risk by
\[
\mathcal P_\gamma^{\{0,B\}}
:=
\left\{
P\in\mathcal P_\gamma:
P\text{ admits a causal completion with }
Y(0),Y(1)\in\{0,B\}\text{ almost surely}
\right\},
\]
\[
R_{\mathrm{bd}}(n,\gamma,x_0)
:=
\inf_T\sup_{P\in\mathcal P_\gamma^{\{0,B\}}}
\mathbb E_{P^{\otimes n}}
|T(\mathcal D_n,x_0)-\mu_{1,P}(x_0)|.
\]
The infimum is over sample-measurable scalar estimators.

Apply \cref{obj:lem:bounded-lower-pair}. It supplies fixed \(a,c,c_0>0\) such that, for every \(n\ge1\), the pair from \cref{obj:def:bounded-lower-pair}, with
\[
h:=ch_{\gamma,n},
\]
belongs to \(\mathcal P_\gamma^{\{0,B\}}\) and satisfies
\[
P_{0,h}\ll P_{1,h},
\qquad
\log\frac{dP_{0,h}}{dP_{1,h}}\in L^1(P_{0,h}),
\qquad
n\,\operatorname{KL}(P_{0,h},P_{1,h})\le\frac18.
\]
Define
\[
\mathbb P_{j,n}:=P_{j,h}^{\otimes n},
\qquad
\tau_{j,n}:=\mu_{1,P_{j,h}}(x_0)
\quad(j\in\{0,1\}),
\qquad
\delta_n:=c_0r_{\gamma,n}.
\]
The separation supplied by \cref{obj:lem:bounded-lower-pair} is
\[
0<\delta_n\le|\tau_{1,n}-\tau_{0,n}|.
\]

For an arbitrary sample-measurable estimator \(T\), define its clipped version by
\[
T_{\mathrm{clip}}(\mathcal D_n)
:=
\operatorname{clip}_{[-B,B]}
\bigl(T(\mathcal D_n,x_0)\bigr).
\]
The response bound established above gives, for \(j\in\{0,1\}\),
\[
|\tau_{j,n}|\le B,
\qquad
|T_{\mathrm{clip}}-\tau_{j,n}|
\le |T(\mathcal D_n,x_0)-\tau_{j,n}|,
\qquad
|T_{\mathrm{clip}}-\tau_{j,n}|\le2B.
\]
Thus the clipped losses are integrable under both sample laws.

Product entropy tensorization and Pinsker's inequality give
\[
\operatorname{KL}(\mathbb P_{0,n},\mathbb P_{1,n})
\le n\,\operatorname{KL}(P_{0,h},P_{1,h})
\le\frac18\le\frac12,
\]
\[
\|\mathbb P_{0,n}-\mathbb P_{1,n}\|_{\mathrm{TV}}
\le
\sqrt{\frac{\operatorname{KL}(\mathbb P_{0,n},\mathbb P_{1,n})}{2}}
\le\frac12.
\]
Define the two large-error events by
\[
\mathcal A_{j,n}
:=
\left\{
\mathcal D_n:
|T_{\mathrm{clip}}(\mathcal D_n)-\tau_{j,n}|
\ge\frac{\delta_n}{2}
\right\},
\qquad j\in\{0,1\}.
\]
The triangle inequality and target separation imply
\[
\mathcal A_{0,n}^c\subseteq\mathcal A_{1,n}.
\]
Hence
\[
\begin{aligned}
\mathbb P_{0,n}(\mathcal A_{0,n})
+\mathbb P_{1,n}(\mathcal A_{1,n})
&\ge
1-\|\mathbb P_{0,n}-\mathbb P_{1,n}\|_{\mathrm{TV}},
\\
\max_{j\in\{0,1\}}\mathbb P_{j,n}(\mathcal A_{j,n})
&\ge
\frac{1-
\sqrt{\operatorname{KL}(\mathbb P_{0,n},\mathbb P_{1,n})/2}}{2}
\ge\frac14.
\end{aligned}
\]
For each \(j\), integration over \(\mathcal A_{j,n}\) yields
\[
\mathbb E_{\mathbb P_{j,n}}
|T_{\mathrm{clip}}-\tau_{j,n}|
\ge\frac{\delta_n}{2}\mathbb P_{j,n}(\mathcal A_{j,n}).
\]
Combining the two inequalities and the clipping contraction gives
\[
\begin{aligned}
\max_{j\in\{0,1\}}
\mathbb E_{\mathbb P_{j,n}}
|T(\mathcal D_n,x_0)-\tau_{j,n}|
&\ge
\max_{j\in\{0,1\}}
\mathbb E_{\mathbb P_{j,n}}
|T_{\mathrm{clip}}-\tau_{j,n}|
\\
&\ge\frac{\delta_n}{8}
=\frac{c_0}{8}r_{\gamma,n}.
\end{aligned}
\]
Both laws belong to the restricted subclass. Taking the infimum over \(T\), and defining
\[
c_\gamma:=\frac{c_0}{8}>0,
\]
proves
\[
c_\gamma r_{\gamma,n}
\le R_{\mathrm{bd}}(n,\gamma,x_0)
\qquad(n\ge1).
\]
% lean: CausalSmith.Stat.WeakOverlap.boundedPointwiseRisk_oracle_lower

\emph{5. Constants for all positive sample sizes.}
For the fixed \(\gamma\), define a local adaptation interval by
\[
\gamma_{\mathrm{lo}}:=\frac{\gamma+1}{2},
\qquad
\gamma_{\mathrm{hi}}:=\gamma+1.
\]
Since \(\gamma>1\),
\[
1<\gamma_{\mathrm{lo}}<\gamma_{\mathrm{hi}},
\qquad
\gamma\in[\gamma_{\mathrm{lo}},\gamma_{\mathrm{hi}}].
\]
Applying the uniform estimator bound to this interval supplies \(K_{\mathrm{loc}}>0\) and \(n_{\mathrm{loc}}\in\mathbb N\) such that
\[
\sup_{P\in\mathcal P_\gamma}
\mathbb E_{P^{\otimes n}}
\left[
\sup_{x\in\mathcal X}
|\widehat\mu_n(x)-\mu_{1,P}(x)|
\right]
\le K_{\mathrm{loc}}r_{\gamma,n}
\quad(n\ge\max\{n_{\mathrm{loc}},1\}).
\]
Define
\[
N_\gamma:=\max\{n_{\mathrm{loc}},1\},
\qquad
q_{\mathrm{small}}:=N_\gamma^{-2},
\qquad
K_{\mathrm{small}}:=\frac{2B}{q_{\mathrm{small}}},
\]
\[
K_\gamma
:=
\max\{c_\gamma,\max\{K_{\mathrm{loc}},K_{\mathrm{small}}\}\}.
\]
These constants are finite, \(q_{\mathrm{small}}>0\), \(K_{\mathrm{small}}>0\), and
\[
0<c_\gamma\le K_\gamma,
\qquad
K_{\mathrm{loc}}\le K_\gamma,
\qquad
K_{\mathrm{small}}\le K_\gamma.
\]

Fix \(n\ge1\). If \(n\ge n_{\mathrm{loc}}\), the measurable estimator \(\widehat\mu_n\) is an admissible witness in the supremum minimax infimum, so
\[
R_\infty(n,\gamma)
\le
\sup_{P\in\mathcal P_\gamma}
\mathbb E_{P^{\otimes n}}
\left[
\sup_{x\in\mathcal X}
|\widehat\mu_n(x)-\mu_{1,P}(x)|
\right]
\le K_{\mathrm{loc}}r_{\gamma,n}
\le K_\gamma r_{\gamma,n}.
\]
If \(n<n_{\mathrm{loc}}\), then \(1\le n\le N_\gamma\). Monotonicity of the inverse square and \(n^{-2}\le r_{\gamma,n}\) give
\[
q_{\mathrm{small}}
=N_\gamma^{-2}
\le n^{-2}
\le r_{\gamma,n}.
\]
Therefore
\[
2B
=K_{\mathrm{small}}q_{\mathrm{small}}
\le K_{\mathrm{small}}r_{\gamma,n}
\le K_\gamma r_{\gamma,n}.
\]
The same clipped estimator has loss at most \(2B\) for every sample, so
\[
R_\infty(n,\gamma)
\le
\sup_{P\in\mathcal P_\gamma}
\mathbb E_{P^{\otimes n}}
\left[
\sup_{x\in\mathcal X}
|\widehat\mu_n(x)-\mu_{1,P}(x)|
\right]
\le2B
\le K_\gamma r_{\gamma,n}.
\]
Thus the upper bound holds for every \(n\ge1\).
% lean: CausalSmith.Stat.WeakOverlap.adaptive_minimax_rate

\emph{6. Comparison of the risks.}
The inclusion
\[
\mathcal P_\gamma^{\{0,B\}}\subseteq\mathcal P_\gamma
\]
implies, for each scalar estimator \(T\),
\[
\sup_{P\in\mathcal P_\gamma^{\{0,B\}}}
\mathbb E_{P^{\otimes n}}
|T(\mathcal D_n,x_0)-\mu_{1,P}(x_0)|
\le
\sup_{P\in\mathcal P_\gamma}
\mathbb E_{P^{\otimes n}}
|T(\mathcal D_n,x_0)-\mu_{1,P}(x_0)|.
\]
Taking infima over the same estimators gives
\[
R_{\mathrm{bd}}(n,\gamma,x_0)
\le R_{\mathrm{pt}}(n,\gamma,x_0).
\]
For any admissible function-valued estimator \(T\), define its evaluation by
\[
T_{x_0}(\mathcal D_n):=T(\mathcal D_n)(x_0).
\]
This is a measurable scalar estimator, and \(x_0\in\mathcal X\) gives
\[
|T_{x_0}(\mathcal D_n)-\mu_{1,P}(x_0)|
\le
\sup_{x\in\mathcal X}
|T(\mathcal D_n)(x)-\mu_{1,P}(x)|.
\]
Taking expectations, worst-case risks, and then the infimum over function-valued estimators yields
\[
R_{\mathrm{pt}}(n,\gamma,x_0)\le R_\infty(n,\gamma).
\]
Combining the bounded-subclass lower bound, the upper bound for every positive sample size, and these comparisons proves
\[
c_\gamma r_{\gamma,n}
\le R_{\mathrm{bd}}(n,\gamma,x_0)
\le R_{\mathrm{pt}}(n,\gamma,x_0)
\le R_\infty(n,\gamma)
\le K_\gamma r_{\gamma,n},
\qquad n\ge1.
\]
% lean: CausalSmith.Stat.WeakOverlap.boundedPointwiseRisk_le_pointwiseRisk
% lean: CausalSmith.Stat.WeakOverlap.pointwiseRisk_le_supremumRisk
\end{proof}
